% !TEX program = xelatex
\documentclass[UTF8,fontset=fandol,a4paper,11pt]{ctexart}
\usepackage{../../_manual_common/hysim_manual}

\title{\textbf{交直流混合高弹性能源电力系统仿真平台}\\[0.5em]
       {\Large 园区综合能源系统优化技术专著}\\[0.3em]
       {\large 电--热--氢--燃料耦合、分层储能、交通与碳约束}}
\author{HySim-XJTU-HRPES（西安交通大学 HRPES 团队）}
\date{\today}

\begin{document}
\maketitle
\begin{abstract}
本文档对应 \srcpath{src/integrated\_energy/} 与
\srcpath{include/hacdcpf/integrated\_energy/}，以园区多载能线性/混合整数优化为对象，
完整说明数据归一化、电--热--氢--燃料平衡、五类库存、分层储氢、交通能源需求、
CCUS、四类目标、后端选择和结果归因。手册采用“理论--实现--接口--数值证据--准入”
双层结构：理论章给出标准能源枢纽模型及成立条件，源码等价章只转写当前 C++ 路径。

当前求解器不接收 \texttt{HybridPowerSystem}，\fld{pcc\_ac\_bus} 只是稳定归因标签；
模型不包含交流无功、电压、支路潮流或安全约束。任何 \fld{feasible} 结果因此只证明
所建多载能代数模型存在后端原始解，不构成电网可行性证书。
\end{abstract}

\tableofcontents
\newpage

\section{范围、架构与阅读契约}

\subsection{模块定位}

公共入口只有
\begin{equation}
R=\operatorname{solve\_campus\_ies}(D,O),
\end{equation}
其中 $D$ 为 \texttt{CampusIESData}，$O$ 为 \texttt{CampusIESOptions}，$R$ 为
\texttt{CampusIESResult}。实现锚点为
\srcpath{src/integrated\_energy/integrated\_energy\_optimizer.cpp:solve\_campus\_ies}。
函数先复制并归一化输入，再用 MIPSolvers AML 创建变量、约束和目标，调用所选后端，最后把
原始解提取为逐时序列、汇总指标和 Sankey 边。函数没有网络模型参数，也不调用 PF/OPF。

\begin{center}\footnotesize
\begin{tabular}{@{}L{3.0cm}L{4.8cm}L{7.0cm}@{}}
\toprule
层 & 公共类型/实现 & 职责\\
\midrule
输入 & \texttt{CampusIESData} & 时域、需求、可用功率、容量、效率、价格、交通、碳与 CCUS 数据\\
选项 & \texttt{CampusIESOptions} & 后端、目标、互斥开关、预算、末态和求解容差\\
建模 & \texttt{solve\_campus\_ies} & AML 变量、线性等式/不等式、可选二进制变量和目标\\
结果 & \texttt{CampusIESResult} & 求解状态、$n$ 项功率/流量、$n+1$ 项库存、总量与有效性声明\\
示例 & \texttt{make\_sample\_campus\_ies\_data} & 生成可运行的合成 24 h 风光荷价曲线，不是标准算例\\
\bottomrule
\end{tabular}
\end{center}

\subsection{载能、符号与单位}

功率/流率变量统一以 MW 表示，库存以 MWh 表示，时间步长以 h 表示；因此
$\Delta t P_t$ 为 MWh。氢与燃料也使用能量等价值 MW/MWh，而不是 kg、Nm$^3$ 或低位热值质量流率。
碳排放和捕集量为每步 tCO$_2$；排放因子为 tCO$_2$/MWh；交通需求为 km，
车型能耗强度为 MWh/km。成本字段沿其乘数决定量纲，最终 \fld{total\_cost} 是调用者价格单位。

电网交换采用两套符号：优化变量 $P_t^{imp},P_t^{exp}\ge0$；结果
\begin{equation}
P_t^{pcc}=P_t^{exp}-P_t^{imp},
\end{equation}
故正值表示园区向外送电。该符号由
\srcpath{src/integrated\_energy/integrated\_energy\_optimizer.cpp:solve\_campus\_ies}
提取循环定义。

\subsection{时域与索引空间}

所有逐时决策向量按输入时间位置 $t=0,\ldots,n-1$ 排列，不使用稳定组件 ID。
五类库存向量按状态节点 $t=0,\ldots,n$ 排列，首项是输入初态，末项是时域结束状态。
\fld{pcc\_ac\_bus} 是唯一稳定外部标签；它不用于矩阵索引和电气查找。

\subsection{能力边界}

\begin{itemize}
  \item 模型是集中式园区能源枢纽，不表达载能网络节点、管道压降、热网温降或氢气压力；
  \item PCC 只有聚合有功交换，无功被明确排除，\fld{fixed\_power\_factor} 非 1 时拒绝；
  \item 风光出力为可用量拆分，未建模随机性、爬坡、最小开停机或预测误差；
  \item CHP 是凸可行域上界，不是固定热电比，也没有启停二进制；
  \item EV/HV/ICV 比例由输入归一化后固定，当前 $P^{EV,V2G}$ 上下界均为零；
  \item 互斥开关开启时才形成 MILP；全部关闭时当前模型为 LP；
  \item Sankey 边由解后汇总生成，不是独立约束，也不提供损耗闭合证明。
\end{itemize}

\subsection{推荐阅读顺序}

研究读者先读第 2 章的标准枢纽模型，再用第 3 章确认当前代码采用的简化；接口调用者重点阅读
第 4--5 章的默认值、长度、失败和索引契约；结果审核者应同时使用第 6 章的残差复算与准入表。


% theory_multi_carrier_optimization.tex —— 综合能源通用理论层
% 本文件遵循 docs/modules/THEORY_WRITING_GUIDE.md。
\section{多载能系统优化理论}

\begin{theorynote}
本章为通用数学理论，给出离散时间能源枢纽、转换矩阵、库存递推、碳核算和多目标优化的
统一表述。公式不要求逐式对应代码；当前实现的对应和缺口在章末明确列出。
\end{theorynote}

\subsection{载能网络与转换矩阵}

令载能集合 $\mathcal C=\{e,q,h,f\}$ 分别表示电、热、氢和燃料。在时段 $t$，
外部输入、内部转换、储能交换和终端需求满足
\begin{equation}
\mathbf u_t+\mathbf C\mathbf x_t+\mathbf d_t^{dis}
=\boldsymbol\ell_t+\mathbf d_t^{ch}+\mathbf w_t.
\label{eq:ies-hub-balance}
\end{equation}
其中 $\mathbf C$ 的每一列描述一个转换设备对各载能母线的有符号注入，$\mathbf w_t$
包含弃能、辅助用能和损耗。式~\eqref{eq:ies-hub-balance} 是能源枢纽模型：它保证每个载能
逐时守恒，但不表达载能输送网络的空间状态。

若设备 $j$ 以载能 $a$ 输入、以 $b$ 输出，固定效率模型为
\begin{equation}
x_{bjt}=\eta_jx_{ajt},\qquad 0<\eta_j\le1.
\end{equation}
热泵的 COP 是环境热源参与下的性能系数，可大于 1，不能与被动能量转换效率混用。

\subsection{CHP 凸可行域}

多输出 CHP 的常用线性外逼近可写为
\begin{equation}
0\le P_t\le\eta_e F_t,\quad
0\le Q_t\le\eta_h F_t,\quad
P_t+Q_t\le\eta_{tot}F_t.
\label{eq:ies-chp-envelope}
\end{equation}
它定义以燃料输入为尺度的凸多面体。若 $\eta_{tot}\le\eta_e+\eta_h$，第三个平面限制
同时满发；若更大则可能冗余。该模型允许在可行域内部自由选择热电比，不能解释为背压机固定耦合曲线。

\begin{gapnote}
非凸 CHP 启停、最小负荷、抽凝机分段可行域、爬坡与启停成本属于机组组合模型；当前园区实现
没有这些变量和约束。
\end{gapnote}

\subsection{库存递推与能量守恒}

对任一储能 $s$，状态递推为
\begin{equation}
E_{s,t+1}=\rho_sE_{s,t}+\eta_s^{ch}\Delta t P_{s,t}^{ch}
-\frac{\Delta t}{\eta_s^{dis}}P_{s,t}^{dis}+\Delta t\sum_r
\left(\eta_{rs}F_{rs,t}-F_{sr,t}\right).
\label{eq:ies-storage}
\end{equation}
其中 $0<\rho_s\le1$ 是每步保持率。若输入保持率本意为“每小时”，步长不为 1 h 时应使用
$\rho_s^{\Delta t}$；若本意为“每步”，则直接使用 $\rho_s$。二者必须在接口契约中区分。

循环末态 $E_{s,T}=E_{s,0}$ 消除时域末端免费放空/充满效应，但也意味着研究窗口代表周期运行。
对季节储氢，短时域强制循环可能压制合理的跨季净转移，必须结合时域解释。

\begin{proposition}[无互斥变量时的同时充放]
若充放均有严格正吞吐成本、$\eta^{ch}\eta^{dis}<1$，且不存在负价格或其他奖励抵消损耗，
则任一最优解都存在一个等价或更优解，使同一储能同一时段不同时充放。
\end{proposition}
\begin{proof}
若 $P^{ch},P^{dis}>0$，沿保持净状态变化不变的方向同时减少二者，直至至少一个为零。
载能母线净需求因往返损耗减少而不增加，严格正吞吐成本下降，故原点不可能是唯一更优结构。
当电价为负、弃能惩罚或跨层奖励存在时，此论证的目标单调性可能失效，需显式互斥。
\end{proof}

\subsection{交通能源映射}

若总里程 $D_t$ 和车型比例 $r_m$ 外生，则
\begin{equation}
D_{m,t}=r_mD_t,\qquad
P_{m,t}=\frac{\alpha_mD_{m,t}}{\Delta t},\qquad \sum_mr_m=1.
\end{equation}
这是需求映射，不是交通分配优化。若 $r_m$ 成为决策变量，则需要车辆拥有量、补能能力、出行服务
和基础设施容量等约束，否则会把全部里程移向目标最便宜的载能。

\subsection{碳流、捕集和残余排放}

外购电和燃料的直接核算可写为
\begin{equation}
E_t^{gross}=g_t^e\Delta tP_t^{imp}+g^f\Delta tF_t^{external}.
\end{equation}
捕集量满足 $0\le C_t\le\phi E_t^{gross}$，残余排放物理等式应为
$R_t=E_t^{gross}-C_t$。若模型只施加 $R_t\ge E_t^{gross}-C_t$，则必须由正碳成本、碳目标或预算
使其取紧；否则 $R_t$ 不是唯一物理量。

\begin{gapnote}
当前实现核算运行期直接排放，不含设备制造、燃料上游、网络损耗归因、边际排放因子反馈或碳库存。
\end{gapnote}

\subsection{多目标标度}

加权和目标
\begin{equation}
J=w_C C+w_E E+w_U U
\end{equation}
只有在权重吸收三项量纲和数量级后才有明确偏好含义。词典序近似
$E+\epsilon C$ 要求 $\epsilon$ 足够小，使任意可辨的 $E$ 改善都优先于 $C$，同时又大于求解器数值噪声。
固定 $10^{-6}$ 只是实现选择，不对任意价格规模自动保证严格词典序。

\subsection{与实现的对应}

\begin{center}\footnotesize
\begin{longtable}{@{}L{7.0cm}L{8.2cm}@{}}
\toprule
理论条目 & 实现状态\\\midrule
\endfirsthead
\toprule 理论条目 & 实现状态\\\midrule\endhead
四载能逐时平衡（式~\eqref{eq:ies-hub-balance}） &
\implfull{src/integrated\_energy/integrated\_energy\_optimizer.cpp:solve\_campus\_ies}\\
CHP 凸包络（式~\eqref{eq:ies-chp-envelope}） &
\implfull{src/integrated\_energy/integrated\_energy\_optimizer.cpp:solve\_campus\_ies}\\
五类库存与四条层间转移（式~\eqref{eq:ies-storage}） &
\implpart{保持率按每步直接乘，不按 $\Delta t$ 指数化；层间方向不互斥}\\
循环末态 & \implfull{src/integrated\_energy/integrated\_energy\_optimizer.cpp:solve\_campus\_ies}\\
交通需求映射 & \implpart{比例固定并归一化；V2G 变量固定为零，车型比例不优化}\\
碳捕集物理等式 & \implpart{残余碳只有下界；需目标或预算保证取紧}\\
运行碳核算 & \implfull{src/integrated\_energy/integrated\_energy\_optimizer.cpp:solve\_campus\_ies}\\
载能网络状态与安全约束 & \implnone\\
随机/鲁棒多阶段优化 & \implnone\\
\bottomrule
\end{longtable}
\end{center}

\subsection{参考文献}
{\renewcommand{\section}[2]{}%
\begin{thebibliography}{9}\small
\bibitem{ies:geidl2007} M.~Geidl, G.~Koeppel, P.~Favre-Perrod, B.~Klockl,
G.~Andersson, and K.~Frohlich, Energy hubs for the future,
\emph{IEEE Power and Energy Magazine}, vol.~5, no.~1, pp.~24--30, 2007.
\bibitem{ies:morvaj2016} B.~Morvaj, R.~Evins, and J.~Carmeliet,
Optimising urban energy systems: simultaneous system sizing, operation and district heating network layout,
\emph{Energy}, vol.~116, pp.~619--636, 2016.
\bibitem{ies:conejo2010} A.~J. Conejo, M.~Carrion, and J.~M. Morales,
\emph{Decision Making Under Uncertainty in Electricity Markets}, Springer, 2010.
\end{thebibliography}}


% theory_conservation_storage.tex -- 综合能源守恒、库存与互斥理论
% 本文件遵循 docs/modules/THEORY_WRITING_GUIDE.md。
\section{载能守恒、库存动力学与无套利条件}

\begin{theorynote}
本章从载能图的关联矩阵出发，推导逐时守恒、跨时域库存恒等式以及导入/导出与充/放互斥的
充分条件。结论用于解释数值证据应检查哪些残差；只有章末对应表标为“完整实现”的条目才是当前
园区求解器的运行契约。
\end{theorynote}

\subsection{载能图与符号}

令载能集合为 $\mathcal C=\{e,q,h,f\}$，时段集合为
$\mathcal T=\{0,\ldots,T-1\}$。把每个外部源、负荷、转换器、库存端口写成有向超边，
$A\in\mathbb R^{|\mathcal C|\times m}$ 为载能--端口关联矩阵：向载能母线注入取正，取出取负。
设备活动向量 $x_t\in\mathbb R^m$ 与外生净需求 $d_t$ 满足
\begin{equation}
  A x_t=d_t,\qquad t\in\mathcal T .
  \label{eq:ies-incidence-balance}
\end{equation}
它是四条逐时平衡的紧凑写法。转换效率不应塞进单位关联矩阵，而应由设备局部约束
$B_jx_{jt}=0$ 表达；这样可以区分“端口拓扑”与“能量损耗”。

\begin{proposition}[逐时守恒的全时域恒等式]
若式~\eqref{eq:ies-incidence-balance} 对所有 $t$ 成立，则任意载能 $c$ 的累计外部输入等于
累计终端需求、净库存吸收、转换净取出和弃能之和。若所有量以 MW 表示，累计式必须乘
$\Delta t$ 后才具有 MWh 量纲。
\end{proposition}
\begin{proof}
取 $A$ 的第 $c$ 行，对 $t$ 求和并乘 $\Delta t$。有限和保持等式；再按端口类型对列分组，
即得所述累计恒等式。该证明不依赖目标函数或求解算法。
\end{proof}

该命题给出两级残差：逐时残差
\begin{equation}
 r_{c,\infty}=\max_t\left|(Ax_t-d_t)_c\right|
 \label{eq:ies-carrier-residual}
\end{equation}
用于发现某一时段的符号或单位错误；累计残差可能因正负抵消而很小，不能替代
式~\eqref{eq:ies-carrier-residual}。

\subsection{单库存递推的望远镜恒等式}

库存 $s$ 的每步模型写为
\begin{equation}
 E_{t+1}=\rho E_t+\eta^{ch}\Delta t P_t^{ch}
 -\frac{\Delta t}{\eta^{dis}}P_t^{dis}+\Delta t\,\xi_t,
 \label{eq:ies-stock-recursion-expanded}
\end{equation}
其中 $\xi_t$ 是层间净输入。反复代入可得闭式解
\begin{equation}
 E_T=\rho^T E_0+\sum_{t=0}^{T-1}\rho^{T-1-t}\Delta t
 \left(\eta^{ch}P_t^{ch}-\frac{P_t^{dis}}{\eta^{dis}}+\xi_t\right).
 \label{eq:ies-stock-closed-form}
\end{equation}

\begin{proposition}[循环边界下的损耗恒等式]
若 $E_T=E_0$，则
\begin{equation}
 \sum_t\rho^{T-1-t}\Delta t
 \left(\eta^{ch}P_t^{ch}-P_t^{dis}/\eta^{dis}+\xi_t\right)
 =(1-\rho^T)E_0 .
 \label{eq:ies-cyclic-loss-identity}
\end{equation}
当 $\rho=1$ 且无层间转移时，累计放电的库存侧能量不能超过累计充电的库存侧能量；循环条件下二者相等。
\end{proposition}
\begin{proof}
在式~\eqref{eq:ies-stock-closed-form} 中代入 $E_T=E_0$ 并移项即可。右端是为补偿
静置损耗所需的折算净输入。
\end{proof}

当前接口中的 $\rho$ 是“每步保持率”。若物理数据给的是每小时保持率 $\rho_h$，任意步长下应使用
$\rho=\rho_h^{\Delta t}$；直接使用 $\rho_h$ 会使时间离散变化时产生模型漂移。

\subsection{分层储氢的状态空间}

把日、周、季三层状态写成 $z_t=(H_t^d,H_t^w,H_t^s)^\top$，则
\begin{equation}
 z_{t+1}=Rz_t+\Delta t(Bu_t+Gv_t),
 \label{eq:ies-layered-state}
\end{equation}
其中 $R=\operatorname{diag}(\rho_d,\rho_w,\rho_s)$；$u_t$ 为各层充放；$v_t$ 为
日到周、周到日、周到季、季到周四条转移。$G$ 每列在源层取 $-1$，在目的层取对应效率。
因此每一次层间转移都使“未加权总氢量”减少
$(1-\eta_{ab})\Delta t v_{ab,t}$。

\begin{corollary}[层间转移不创造能量]
若所有转移效率位于 $(0,1]$，则对 $\mathbf 1^\top z_t$ 求和，层间流只可能保持或降低三层总库存，
不会凭空增加氢能。
\end{corollary}

\begin{gapnote}
当前实现允许同一时段同时存在相反方向的层间转移；储能互斥二进制只约束各层自身充放，未约束
日--周和周--季的双向流。正吞吐成本通常压制循环流，但这不是结构保证。
\end{gapnote}

\subsection{导入与导出的无套利充分条件}

电力交换满足净注入 $g_t=P_t^{imp}-P_t^{exp}$。固定其他变量，保持 $g_t$ 不变而同时减少
$P_t^{imp},P_t^{exp}$ 的方向是 $(-\delta,-\delta)$。

\begin{proposition}[价格次序排除同时购售]
若购电边际价 $c_t^{buy}$ 严格大于售电边际收入 $c_t^{sell}$，且减少购售不改变其他约束，
则任何最优解都可取 $P_t^{imp}P_t^{exp}=0$。若严格不等式成立，含同时购售的点不可能最优。
\end{proposition}
\begin{proof}
令 $\delta=\min(P_t^{imp},P_t^{exp})>0$。净交换不变，目标变化为
$-\delta(c_t^{buy}-c_t^{sell})<0$，与最优性矛盾。
\end{proof}

负购电价、售电补贴、分段价格或依赖毛流量的碳/网络费会破坏证明条件。因此工程准入不能仅凭“通常
购价高于售价”省略互斥；当前接口提供二进制门显式建立互斥。

\subsection{充放电互斥的精确消元}

设某时段 $P^{ch},P^{dis}>0$。保持库存变化不变的消元方向满足
$\eta^{ch}\delta^{ch}=\delta^{dis}/\eta^{dis}$，即
$\delta^{dis}=\eta^{ch}\eta^{dis}\delta^{ch}$。电母线净需求的变化为
\begin{equation}
 -\delta^{ch}+\delta^{dis}=-(1-\eta^{ch}\eta^{dis})\delta^{ch}\le0.
\end{equation}
若减少净需求总能以非负边际成本消化，且吞吐成本为正，则同时充放不是最优。反之，负价格、强制消纳、
终端奖励或缺少向下调节时，显式二进制互斥仍有必要。

\subsection{可再生比例分母与碳预算下界}

当前类型的比例约束为
\begin{equation}
 \sum_t\Delta t(P_t^{pv}+P_t^{wind})\ge \alpha
 \sum_t\Delta t(P_t^{imp}+P_t^{pv}+P_t^{wind}+P_t^{chp}+P_t^{fc}).
 \label{eq:ies-renewable-mandate}
\end{equation}
右端是所列“供应量”而非终端负荷，也不扣除出口、储能循环和转换损耗。解释政策比例时必须逐项公布
分子与分母，不能只报一个利用率。

对给定负荷与容量，可先解不含碳预算的最小残余碳问题得到 $R_{min}$。预算 $B<R_{min}$ 是不可行的
充分证据；$B\ge R_{min}$ 只消除这一必要障碍，不保证其他载能与容量约束可行。

\subsection{与实现的对应}

\begin{center}\footnotesize
\begin{longtable}{@{}L{7.1cm}L{8.1cm}@{}}
\toprule 理论条目 & 实现状态\\\midrule
\endfirsthead\toprule 理论条目 & 实现状态\\\midrule\endhead
四载能逐时守恒与最大残差 & \implfull{src/integrated\_energy/integrated\_energy\_optimizer.cpp:solve\_campus\_ies}\\
五类库存递推与循环末态 & \implfull{src/integrated\_energy/integrated\_energy\_optimizer.cpp:solve\_campus\_ies}\\
保持率按 $\rho_h^{\Delta t}$ 换算 & \implnone\\
分层储氢状态空间 & \implpart{三层与四条转移已实现；相反方向转移不互斥}\\
购售与各层充放互斥 & \implfull{src/integrated\_energy/integrated\_energy\_optimizer.cpp:solve\_campus\_ies}\\
可再生比例式~\eqref{eq:ies-renewable-mandate} & \implfull{src/integrated\_energy/integrated\_energy\_optimizer.cpp:solve\_campus\_ies}\\
碳预算最小可行下界预检查 & \implnone\\
\bottomrule
\end{longtable}
\end{center}

\subsection{参考文献}
{\renewcommand{\section}[2]{}%
\begin{thebibliography}{9}\small
\bibitem{ies:geidl2007b} M.~Geidl et al., Energy hubs for the future,
\emph{IEEE Power and Energy Magazine}, vol.~5, no.~1, pp.~24--30, 2007.
\bibitem{ies:boyd2004} S.~Boyd and L.~Vandenberghe, \emph{Convex Optimization},
Cambridge University Press, 2004.
\end{thebibliography}}

% theory_optimization_duality.tex -- 综合能源优化、对偶与数值准入理论
% 本文件遵循 docs/modules/THEORY_WRITING_GUIDE.md。
\section{LP/MILP 结构、对偶解释与数值准入}

\begin{theorynote}
本章说明园区模型为何在关闭互斥时是线性规划、开启互斥时成为混合整数线性规划，并给出 LP 强对偶、
影子价格和误差准入的严格解释。对偶价格、冲突精炼与尺度自动化属于标准理论，不代表当前结果结构已经输出。
\end{theorynote}

\subsection{标准形与模型规模}

汇总所有时段连续变量为 $x$、互斥变量为 $z$，模型可写成
\begin{equation}
 \min_{x,z}\ c^\top x+d^\top z,
 \quad Ax=b,\quad Gx+Hz\le h,
 \quad \ell\le x\le u,quad z\in\{0,1\}^{n_z}.
 \label{eq:ies-milp-standard}
\end{equation}
固定 $z$ 或关闭互斥后即为 LP。变量和局部约束数量均为 $O(T)$；库存递推只连接相邻时段，故约束矩阵
具有块双对角骨架。分支定界的最坏复杂度仍随 $n_z$ 指数增长，$O(T)$ 的模型尺寸不意味着 MILP
求解时间线性增长。

\begin{proposition}[连续模型的凸性]
当效率、COP、价格、排放因子和需求均为给定参数，且不引入二元互斥时，可行域是闭凸多面体；若
变量上下界使其有界且可行，则线性目标存在全局最优解。
\end{proposition}
\begin{proof}
等式、线性不等式和盒约束的交集为闭凸多面体。有界非空多面体是紧集，连续线性函数在其上达到最小值。
\end{proof}

\subsection{LP 对偶与载能边际价格}

为简洁起见，将 LP 写为
\begin{equation}
 \min c^\top x\quad\text{s.t.}\quad Ax=b,\ Gx\le h,\ x\ge0.
\end{equation}
其拉格朗日函数
$L=c^\top x+\lambda^\top(Ax-b)+\mu^\top(Gx-h)$，其中 $\mu\ge0$。
对偶为
\begin{equation}
 \max_{\lambda,\mu}\ -b^\top\lambda-h^\top\mu,
 \quad c+A^\top\lambda+G^\top\mu\ge0,\quad \mu\ge0.
 \label{eq:ies-lp-dual}
\end{equation}
在采用一致符号后，电、热、氢、燃料平衡行的 $\lambda_{c,t}$ 可解释为相应需求增加一单位对最优成本的
局部敏感度。该解释要求最优基稳定；在退化点，单个对偶值可能不唯一。

\begin{theorem}[强对偶与最优性证书]
若原 LP 与对偶 LP 至少一方存在有限最优解并满足标准可行条件，则二者最优值相等。一个原始--对偶点
若满足原始可行、对偶可行和互补松弛，即构成全局最优性证书。
\end{theorem}

MILP 的最终 incumbent 不能用式~\eqref{eq:ies-lp-dual} 直接证明全局最优；需要分支定界下界 $L$ 与
incumbent $U$。常见相对 gap 为
\begin{equation}
 g=\frac{|U-L|}{\max(1,|U|)}.
 \label{eq:ies-mip-gap}
\end{equation}
不同后端可能使用不同分母或绝对 gap 组合，手册只能按适配层实际返回的状态解释。

\subsection{互斥二进制的 Big-M 强度}

购售互斥可写为
\begin{equation}
 0\le P_t^{imp}\le \bar P^{imp}z_t,\qquad
 0\le P_t^{exp}\le \bar P^{exp}(1-z_t).
 \label{eq:ies-grid-binary}
\end{equation}
这里 $M$ 不是任意大数，而应直接取物理交换上界。过大的 $M$ 会削弱 LP 松弛、恶化数值尺度；小于真实
上界则错误裁剪可行域。储能互斥同理，应以充放功率铭牌上界构造。

\subsection{多目标量纲与 Pareto 前沿}

成本 $C$、排放 $R$ 与弃能 $U$ 的单位分别可能为货币、tCO$_2$ 与 MWh。直接写
$w_CC+w_RR+w_UU$ 时，权重带有逆单位。可采用锚点归一化
\begin{equation}
 \hat C=\frac{C-C^{min}}{C^{ref}-C^{min}},\quad
 \hat R=\frac{R-R^{min}}{R^{ref}-R^{min}},\quad
 \hat U=\frac{U-U^{min}}{U^{ref}-U^{min}},
 \label{eq:ies-objective-normalization}
\end{equation}
再用无量纲权重。分母为零时该目标在当前可行域无变化，应移除而不是除零。

当前 Carbon 与 MinCurtailment 目标用 $10^{-6}$ 成本作 tie-break。严格词典序要求次级目标的最大变化乘
系数小于主目标的最小可分辨改善；固定系数对任意规模不自动满足此条件。

\begin{gapnote}
当前实现不生成 Pareto 前沿、不自动归一化权重，也不输出平衡约束对偶变量。手册中的对偶理论用于说明
应如何扩展验证，不得把结果中的成本向量解释为节点或载能边际价格。
\end{gapnote}

\subsection{后端状态与准入谓词}

研究级准入应是合取而不是只读一个布尔量：
\begin{equation}
 \mathcal A=(\text{has primal})\land(\text{termination acceptable})\land
 (r_p\le\tau_p)\land(g\le\tau_g)\land(\text{scope sufficient}).
 \label{eq:ies-admission}
\end{equation}
其中 $r_p$ 至少含四载能平衡、五类库存递推、转换等式、可再生可用量、交通映射与碳关系的最大残差。
若研究问题要求电压或支路安全，当前 aggregate-PCC scope 永远不满足最后一项，无论 MILP 多么精确。

\subsection{数值尺度与误差预算}

MW/MWh 变量通常在 $10^{-3}$--$10^2$，交通里程可达 $10^2$--$10^5$，成本系数可跨多个数量级。
可用对角尺度 $D_x,D_r$ 将 $x=D_x\hat x$、约束左乘 $D_r$。尺度变换不改变物理解，却会改变后端的绝对
可行性容差所代表的物理误差。报告残差时应恢复到 MW、MWh、tCO$_2$ 等工程单位。

独立重算的误差预算可分为
\begin{equation}
 \epsilon_{obs}\le\epsilon_{solver}+\epsilon_{extract}+\epsilon_{oracle},
\end{equation}
其中求解器容差、结果抽取舍入和 oracle 浮点误差来源不同。若误差超过阈值，首先比较逐式残差定位，
不能仅通过放宽统一阈值使测试变绿。

\subsection{与实现的对应}

\begin{center}\footnotesize
\begin{longtable}{@{}L{7.1cm}L{8.1cm}@{}}
\toprule 理论条目 & 实现状态\\\midrule
\endfirsthead\toprule 理论条目 & 实现状态\\\midrule\endhead
式~\eqref{eq:ies-milp-standard} 的连续/二元切换 & \implfull{src/integrated\_energy/integrated\_energy\_optimizer.cpp:solve\_campus\_ies}\\
物理上界 Big-M & \implfull{src/integrated\_energy/integrated\_energy\_optimizer.cpp:solve\_campus\_ies}\\
四类目标与固定 tie-break & \implfull{src/integrated\_energy/integrated\_energy\_optimizer.cpp:solve\_campus\_ies}\\
对偶载能价格与 KKT 输出 & \implnone\\
自动目标/行列尺度化 & \implnone\\
后端 primal、termination 与 gap 归一 & \implpart{结果暴露 primal/optimal/status；不输出完整原始--对偶证书}\\
独立工程单位残差 oracle & \implfull{tools/integrated\_energy\_validation/run\_cross\_validation.py:validate\_case}\\
电网电压与支路安全准入 & \implnone\\
\bottomrule
\end{longtable}
\end{center}

\subsection{参考文献}
{\renewcommand{\section}[2]{}%
\begin{thebibliography}{9}\small
\bibitem{ies:bertsimas1997} D.~Bertsimas and J.~Tsitsiklis,
\emph{Introduction to Linear Optimization}, Athena Scientific, 1997.
\bibitem{ies:nemhauser1988} G.~Nemhauser and L.~Wolsey,
\emph{Integer and Combinatorial Optimization}, Wiley, 1988.
\end{thebibliography}}

\section{源码等价 MILP}

本章逐式对应 \srcpath{src/integrated\_energy/integrated\_energy\_optimizer.cpp:405--777}。
所有功率变量均为非负连续变量；库存变量含 $t=0,\ldots,n$，其余逐时变量含
$t=0,\ldots,n-1$。未在下文展开的上界直接取同名输入字段；$P_{EV,V2G}$ 的上下界均为零，
因此当前实现不允许 V2G（\srcpath{integrated\_energy\_optimizer.cpp:491--503}）。

\subsection{输入归一化}

步数必须为正，步长必须有限且大于零；固定功率因数必须在 $10^{-12}$ 内等于 1，原因是模型
不含无功、电压或电网约束。非有限的一般非负参数通过
$\max(0,x)$ 归零；以下效率与保持率字段经 \texttt{require\_unit\_efficiency} 校验，
必须有限且在 $(0,1]$，否则抛异常：\fld{eta\_electrolysis}、\fld{eta\_power\_to\_fuel}、
\fld{eta\_fuelcell}、\fld{eta\_chp\_elec}、\fld{eta\_chp\_heat}、\fld{eta\_chp\_total}、
\fld{eta\_storage\_charge}、\fld{eta\_storage\_discharge}、五类储能保持率
（\fld{electric\_}/\fld{thermal\_}/\fld{hydrogen\_}/\fld{weekly\_hydrogen\_}/%
\fld{seasonal\_hydrogen\_storage\_retention}）与四个层间转移效率
（\fld{eta\_daily\_to\_weekly}、\fld{eta\_weekly\_to\_daily}、
\fld{eta\_weekly\_to\_seasonal}、\fld{eta\_seasonal\_to\_weekly}）
（\srcpath{integrated\_energy\_optimizer.cpp:164--194}）。\fld{eta\_wasteheat} 与
\fld{eta\_carbon\_to\_fuel} 不参与该校验也不做归一化，按原值直接进入模型
（默认均为 0.0；分别见 \srcpath{integrated\_energy\_optimizer.cpp:564} 与
\srcpath{integrated\_energy\_optimizer.cpp:663--664}）。热泵 COP 非有限或
非正时改为 3.2。EV/HV/ICV 比例先取非负，和大于 $10^{-8}$ 时归一化为和为 1，否则设置为
$(0,0,1)$。每条 profile 长度只能为 0、1 或 $n$；空 profile 使用源码字段对应的固定回退值，
单值扩展为 $n$ 项。依据为 \srcpath{integrated\_energy\_optimizer.cpp:42--245}。

\subsection{逐时载能平衡}

令 $\Delta t=\fld{step\_duration\_hr}$。新能源可用量严格拆分为
\begin{equation}
 P_t^{solar}+P_t^{solar,curt}=\bar P_t^{solar},\qquad
 P_t^{wind}+P_t^{wind,curt}=\bar P_t^{wind},
\end{equation}
且两个电解功率之和不超过同一电解槽功率上限。
电平衡为
\begin{align}
 &P_t^{imp}-P_t^{exp}+P_t^{solar}+P_t^{wind}+P_t^{chp}+P_t^{fc}
 +P_t^{es,dis}+P_t^{EV,V2G}\nonumber\\
 &\quad=L_t^e+P_t^{ely,H2}+P_t^{ely,fuel}+P_t^{hp}+P_t^{es,ch}
 +P_t^{EV,ch}+P_t^{CCUS}.
\end{align}
热平衡为
\begin{equation}
 Q_t^{chp}+Q_t^{hp}+Q_t^{ts,dis}+\eta_{wasteheat}P_t^{hp}
 =L_t^h+Q_t^{ts,ch}.
\end{equation}
依据为 \srcpath{integrated\_energy\_optimizer.cpp:547--566}。

启用分层储氢时，氢平衡为
\begin{equation}
 H_t^{ely}+H_t^{daily,dis}+H_t^{weekly,dis}+H_t^{seasonal,dis}
 =L_t^{H2}+H_t^{fc}+H_t^{HV}+H_t^{daily,ch}+H_t^{weekly,ch}+H_t^{seasonal,ch}.
\end{equation}
关闭时，周/季层充放和四个层间转移变量分别以等式固定为零，并从上述供需式删除周/季充放项。
燃料平衡始终为
\begin{equation}
 F_t^{external}+F_t^{synthetic}=F_t^{chp}+L_t^f+F_t^{ICV}.
\end{equation}
依据为 \srcpath{integrated\_energy\_optimizer.cpp:568--590}。

\subsection{转换设备约束}

代码创建以下等式或不等式：
\begin{align}
 H_t^{ely}&=\eta_{ely}P_t^{ely,H2},&
 F_t^{synthetic}&=\eta_{p2f}P_t^{ely,fuel},\\
 P_t^{fc}&=\eta_{fc}H_t^{fc},&
 Q_t^{hp}&=COP\,P_t^{hp},\\
 P_t^{chp}&\le\eta_{chp,e}F_t^{chp},&
 Q_t^{chp}&\le\eta_{chp,h}F_t^{chp},\\
 P_t^{chp}+Q_t^{chp}&\le\eta_{chp,total}F_t^{chp}.
\end{align}
CHP 的三条关系均为上界，不是固定热电比等式；见
\srcpath{integrated\_energy\_optimizer.cpp:592--605}。

\subsection{库存演化}

电、热和日内氢库存分别为
\begin{align}
 E_{t+1}&=\rho_eE_t+\eta_{ch}\Delta t P_t^{es,ch}
 -\frac{\Delta t}{\eta_{dis}}P_t^{es,dis},\\
 Q_{t+1}^{store}&=\rho_qQ_t^{store}+\eta_{ch}\Delta t Q_t^{ts,ch}
 -\frac{\Delta t}{\eta_{dis}}Q_t^{ts,dis},\\
 H_{t+1}^{d}&=\rho_dH_t^d+\eta_{ch}\Delta t H_t^{d,ch}
 -\frac{\Delta t}{\eta_{dis}}H_t^{d,dis}
 +\Delta t\eta_{w2d}H_t^{w2d}-\Delta t H_t^{d2w}.
\end{align}
周、季库存为
\begin{align}
 H_{t+1}^{w}&=\rho_wH_t^w+\eta_{ch}\Delta t H_t^{w,ch}
 -\frac{\Delta t}{\eta_{dis}}H_t^{w,dis}
 +\Delta t\eta_{d2w}H_t^{d2w}+\Delta t\eta_{s2w}H_t^{s2w}
 -\Delta t H_t^{w2d}-\Delta t H_t^{w2s},\\
 H_{t+1}^{s}&=\rho_sH_t^s+\eta_{ch}\Delta t H_t^{s,ch}
 -\frac{\Delta t}{\eta_{dis}}H_t^{s,dis}
 +\Delta t\eta_{w2s}H_t^{w2s}-\Delta t H_t^{s2w}.
\end{align}
五类库存的初态均通过 fix 固定；启用循环末态时，五类末态也分别固定为各自初始输入。
依据为 \srcpath{integrated\_energy\_optimizer.cpp:505--516} 和
\srcpath{integrated\_energy\_optimizer.cpp:607--639}。

\subsection{交通、碳和 CCUS}

若关闭交通，代码把 $D_t$ 直接设为零；否则使用输入需求。车型里程不是优化分配，而固定为
\begin{align}
 D_t^{EV}&=r_{EV}D_t,&D_t^{HV}&=r_{HV}D_t,&D_t^{ICV}&=r_{ICV}D_t,\\
 P_t^{EV,ch}&=\alpha_{EV}r_{EV}D_t/\Delta t,&
 H_t^{HV}&=\alpha_{HV}r_{HV}D_t/\Delta t,&
 F_t^{ICV}&=\alpha_{ICV}r_{ICV}D_t/\Delta t.
\end{align}
依据为 \srcpath{integrated\_energy\_optimizer.cpp:641--650}。

碳变量按以下约束装配：
\begin{align}
 E_t^{CO2}&=g_t\Delta t P_t^{imp}+g_f\Delta t F_t^{external},\\
 0\le C_t&\le\phi E_t^{CO2},\\
 R_t&\ge E_t^{CO2}-C_t,\\
 P_t^{CCUS}&\ge(c_{power}/\Delta t)C_t.
\end{align}
$C_t$ 另有变量上界 \fld{ccus\_max\_tco2\_per\_h}$\cdot\Delta t$。
若 \fld{eta\_carbon\_to\_fuel}>0，创建
$F_t^{synthetic}\le\eta_{c2f}C_t$；否则固定 $F_t^{synthetic}=0$。
注意 $R_t$ 只有下界而没有等式；其是否取紧由目标或碳预算决定。
依据为 \srcpath{integrated\_energy\_optimizer.cpp:652--668}。

\subsection{互斥、总量约束和目标}

启用电网交换互斥时，以二进制 $z_t$ 创建
\begin{equation}
 P_t^{imp}\le\bar P^{imp}z_t,\qquad P_t^{exp}\le\bar P^{exp}(1-z_t).
\end{equation}
启用储能互斥时，电、热、日/周/季氢五类储能各有独立二进制充电模式，并分别创建
$P^{ch}\le\bar P^{ch}z$、$P^{dis}\le\bar P^{dis}(1-z)$；层间转移没有方向互斥变量。
依据为 \srcpath{integrated\_energy\_optimizer.cpp:518--533} 和
\srcpath{integrated\_energy\_optimizer.cpp:670--708}。

令
\begin{align}
 RENEW&=\sum_t\Delta t(P_t^{solar}+P_t^{wind}),\\
 SUPPLY&=\sum_t\Delta t(P_t^{imp}+P_t^{solar}+P_t^{wind}+P_t^{chp}+P_t^{fc}),\\
 CARBON&=\sum_tR_t,\\
 CURT&=\sum_t\Delta t(P_t^{solar,curt}+P_t^{wind,curt}).
\end{align}
可再生比例大于零时创建 $RENEW\ge\fld{renewable\_mandate\_fraction}\,SUPPLY$；
启用碳预算时创建 $CARBON\le\fld{co2\_budget\_tco2}$。

成本 $COST$ 完全由 \srcpath{integrated\_energy\_optimizer.cpp:716--739} 的线性项组成：购电正、
售电负，外购燃料、新能源/CHP/热泵/电解槽/燃料电池运维、所有储能吞吐、弃能、残余碳、捕集量
和 CCUS 功率均按同名成本字段计价。四个目标分支严格为
\begin{align}
 \texttt{Cost}:&\ COST,\\
 \texttt{Carbon}:&\ CARBON+10^{-6}COST,\\
 \texttt{MinCurtailment}:&\ CURT+10^{-6}COST,\\
 \texttt{Weighted}:&\ w_cCOST+w_eCARBON+w_uCURT.
\end{align}
依据为 \srcpath{integrated\_energy\_optimizer.cpp:710--767}。

\subsection{结果语义}

\fld{feasible} 仅等于 \srcpath{has\_primal()}；\fld{optimal} 还要求后端状态 optimal、gap 有限，
且 $gap\le\fld{mip\_gap\_tol}+10^{-12}$。提取时绝对值小于 $10^{-7}$ 的变量置零。
PCC 报告符号为 $P^{pcc}=P^{export}-P^{import}$，与旧手册中的相反符号概述不同。
依据为 \srcpath{integrated\_energy\_optimizer.cpp:779--849}。

\section{输入数据与选项完整契约}

\subsection{时域与 profile 广播}

\fld{num\_steps} 必须为正，\fld{step\_duration\_hr} 必须有限且大于零。
每个 profile 只接受长度 0、1 或 $n$：空向量用字段专用默认值填充，单值广播到全时域，
长度 $n$ 原样保留，其他长度抛 \texttt{std::invalid\_argument}。实现锚点为
\srcpath{src/integrated\_energy/integrated\_energy\_optimizer.cpp:profile\_or\_default}。

\begin{center}\footnotesize
\begin{longtable}{@{}L{4.3cm}L{2.0cm}L{2.0cm}L{6.3cm}@{}}
\toprule
字段族 & 单位 & 空 profile 回退 & 语义\\\midrule
\endfirsthead
\toprule 字段族 & 单位 & 空 profile 回退 & 语义\\\midrule\endhead
四类 \fld{*\_load\_mw} & MW & 0 & 电、热、氢、燃料终端需求\\
\fld{transport\_demand\_km} & km/步 & 0 & 总出行服务；关闭交通时求解内覆盖为 0\\
\fld{solar/wind\_available\_mw} & MW & 0 & 可用出力，不由 rated 字段自动生成\\
\fld{grid\_buy\_price} & 价格/MWh & 80 & 购电成本\\
\fld{grid\_sell\_price} & 价格/MWh & 35 & 售电收益系数\\
\fld{grid\_carbon\_tco2\_mwh} & tCO$_2$/MWh & 0.55 & 外购电平均排放因子\\
\bottomrule
\end{longtable}
\end{center}

\subsection{容量、库存与效率}

一般容量、功率上限、成本、碳系数和里程强度经 \srcpath{positive\_or\_zero} 处理：
仅有限且大于零的值保留，其余为零。五类库存均有容量、下限、初态、充电上限和放电上限；
初态随后受变量边界检查。周/季氢层与两级双向转移仅在
\fld{enable\_hydrogen\_storage\_layers}=true 时可用。

以下被动效率必须有限且在 $(0,1]$：电解、P2F、燃料电池、CHP 三效率、共同充放效率、
五个保持率和四个层间转移效率。\fld{cop\_heatpump} 非正或非有限时回填 3.2；
\fld{eta\_wasteheat} 与 \fld{eta\_carbon\_to\_fuel} 当前不经过效率校验，按原值进入约束，
这是必须由调用者额外审计的边界。

\subsection{交通与碳字段}

\fld{ev\_ratio}/\fld{hv\_ratio}/\fld{icv\_ratio} 先清洗为非负值；和大于 $10^{-8}$ 时除以总和，
否则替换为 $(0,0,1)$。因此输出车型里程使用归一化比例，不一定等于原始输入。
\fld{alpha\_*\_mwh\_per\_km} 是能耗强度。\fld{fuel\_carbon\_tco2\_mwh} 与电网 profile
共同形成毛排放；\fld{ccus\_capture\_fraction}、\fld{ccus\_max\_tco2\_per\_h} 和
\fld{ccus\_power\_mwh\_per\_tco2} 分别限制捕集比例、捕集率和辅助用电。

\subsection{求解与建模选项}

\begin{paramtable}
\fld{solver} & -- & -- & 枚举 & Auto & AML 后端：Auto、HiGHS、SCIP、Gurobi、Native\\
\fld{objective} & -- & -- & 枚举 & Cost & 成本、碳、最小弃能或加权目标\\
\fld{enable\_grid\_exchange\_exclusivity} & -- & -- & bool & false & 创建逐时购售电方向二进制变量\\
\fld{enable\_storage\_exclusivity} & -- & -- & bool & false & 为五类储能分别创建充放方向二进制变量\\
\fld{enable\_carbon\_budget} & -- & -- & bool & false & 施加全时域残余碳预算\\
\fld{enforce\_terminal\_storage\_cyclic} & -- & -- & bool & true & 五类末态固定为各自初态\\
\fld{enable\_transport} & -- & -- & bool & true & 关闭时交通需求按零处理\\
\fld{enable\_hydrogen\_storage\_layers} & -- & -- & bool & true & 启用周/季库存及层间转移\\
\fld{verbose} & -- & -- & bool & false & AML verbosity 取 1，否则 0\\
\fld{time\_limit\_sec} & -- & s & 正数 & 60 & 透传后端；达到时限仍可能有原始解\\
\fld{mip\_gap\_tol} & $\epsilon_g$ & -- & $\ge0$ & $10^{-4}$ & 后端容差及 \fld{optimal} 二次判定阈值\\
\fld{weight\_cost} & $w_C$ & -- & 实数 & 1 & Weighted 分支成本权重\\
\fld{weight\_carbon} & $w_E$ & -- & 实数 & 0 & Weighted 分支碳权重\\
\fld{weight\_curtailment} & $w_U$ & -- & 实数 & 0 & Weighted 分支弃能权重\\
\end{paramtable}

\fld{solver} 只选择首选后端；\srcpath{solve\_campus\_ies} 设置
\fld{allow\_fallback}=true，故实际 \fld{solver\_name} 可能不同于请求。结果审计必须读取实际后端和状态，
不能用请求枚举推断执行路径。

\subsection{输入拒绝与静默清洗清单}

\begin{center}\footnotesize
\begin{longtable}{@{}L{5.0cm}L{4.5cm}L{5.5cm}@{}}
\toprule
条件 & 行为 & 风险口径\\\midrule
\endfirsthead
\toprule 条件 & 行为 & 风险口径\\\midrule\endhead
$n\le0$ 或 $\Delta t\le0$/非有限 & 抛异常 & 无结果对象\\
profile 长度不是 0、1、$n$ & 抛异常 & 不截断、不循环填充\\
\fld{fixed\_power\_factor}$\ne1$ & 抛异常 & 无功模型未实现\\
被动效率不在 $(0,1]$ & 抛异常 & fail-closed\\
一般容量/成本非正或非有限 & 清零 & 可能把错误输入静默退化为禁用设备\\
COP 非正或非有限 & 回填 3.2 & 必须在外部输入审计记录回退\\
交通比例和近零 & 回填 $(0,0,1)$ & 默认全部视为燃油车\\
\bottomrule
\end{longtable}
\end{center}

\section{求解器状态、结果向量与可视化契约}

\subsection{可行、最优与失败语义}

\fld{feasible} 精确等于 AML \srcpath{SolveResult::has\_primal}。
\fld{optimal} 要求后端报告 optimal、\fld{optimality\_gap} 有限且
\begin{equation}
\fld{optimality\_gap}\le\fld{mip\_gap\_tol}+10^{-12}.
\end{equation}
达到时限或 fallback 后只要存在原始解，\fld{feasible} 可以为真而 \fld{optimal} 为假。
无原始解时函数仍返回状态、实际求解器、目标值和耗时，但不填充时序结果。

\subsection{结果长度与索引}

\begin{center}\footnotesize
\begin{longtable}{@{}L{5.0cm}L{2.0cm}L{8.0cm}@{}}
\toprule
结果族 & 长度 & 单位/口径\\\midrule
\endfirsthead
\toprule 结果族 & 长度 & 单位/口径\\\midrule\endhead
所有 \fld{p\_*}/\fld{q\_*}/\fld{h\_*}/\fld{f\_*} 流率 & $n$ & MW，按输入时间位置\\
\fld{e\_storage\_mwh}、\fld{q\_storage\_mwh}、三层 \fld{h\_*\_storage\_mwh} & $n+1$ & MWh，含初态和末态\\
\fld{d\_ev\_km}/\fld{d\_hv\_km}/\fld{d\_icv\_km} & $n$ & km/步，固定比例映射\\
\fld{emissions\_tco2}/\fld{co2\_captured\_tco2}/\fld{carbon\_residual\_tco2} & $n$ & tCO$_2$/步\\
\fld{p\_pcc\_mw} & $n$ & MW，export-import；\fld{pcc\_ac\_bus} 仅归因标签\\
\bottomrule
\end{longtable}
\end{center}

提取函数把绝对值小于 $10^{-7}$ 的原始变量置为零。这是报告清洗阈值，不是可行性容差；
用清洗后序列复算平衡时应使用与测试一致的正容差，而不能期待位级相等。

\subsection{汇总量}

所有 MW 序列总能量按 $\Delta t\sum_tP_t$ 计算；碳变量已经是每步 tCO$_2$，直接求和。
\fld{total\_transport\_km} 是三类求解里程之和，关闭交通时为零。
\begin{equation}
\fld{renewable\_utilization}=
\begin{cases}
E^{renew}/E^{available},&E^{available}>10^{-9},\\
0,&\text{否则。}
\end{cases}
\end{equation}
\fld{total\_cost} 由提取后序列重新计算，项与成本目标一致；当目标为 Carbon、MinCurtailment 或 Weighted 时，
\fld{objective} 与 \fld{total\_cost} 不相等是预期行为。

\subsection{有效性和模型限制}

\texttt{CampusIESValidity} 的当前默认/返回口径为：多载能平衡和聚合 PCC 有功已建模，
电气网络耦合、无功、电压和支路限值均为 false。\fld{model\_scope} 固定为
\srcpath{isolated-campus-multi-carrier-milp}；\fld{model\_limitations} 至少包含 PCC 未接 AC 网络及
无功/电压/支路/安全约束未建模两项。调用方不得用空 limitations 作为成功条件，也不得在外层删除这些声明。

\subsection{Sankey 边}

\texttt{CampusIESSankeyFlow} 使用英文 source/target、MWh 数值和 carrier 字符串。
\srcpath{add\_sankey\_flow} 只保留大于 $10^{-9}$ 的边，因此零流不会出现。边值由设备流量积分，
包括充放电两条方向边，但没有单列保持损失、转换损失和全部 CCUS 碳流；故 Sankey 是展示视图，
不是可从边集合独立复算的完整守恒图。

\subsection{复杂度与数值风险}

连续变量和逐时约束数量均为 $O(n)$。开启购售互斥增加 $n$ 个二进制；开启储能互斥增加
$5n$ 个二进制，全部开启共 $6n$。LP 路径的实际复杂度取决于后端稀疏线性代数；MILP 最坏复杂度
为指数级，\fld{time\_limit\_sec} 和 gap 是必要证书字段。固定 $10^{-6}$ 目标扰动、$10^{-7}$ 提取清零和
$10^{-9}$ 汇总/Sankey 阈值处在不同语义层，不能互换。

\section{可复现实验、解析解与独立方程交叉验证}

\subsection{证据生成路径}

本章数字由 Release 目标 \texttt{integrated\_energy\_xref} 生成原始 JSON，再由不链接 HySim/MIPSolvers
建模 API 的 Python 程序逐式重算。注册测试命令为
\begin{verbatim}
ctest --test-dir build/macos-release \
  -R '^integrated_energy_cross_validation$' --output-on-failure
\end{verbatim}
生成器位于
\srcpath{tools/integrated\_energy\_validation/validate\_integrated\_energy\_xref.cpp:main}，
oracle 位于
\srcpath{tools/integrated\_energy\_validation/run\_cross\_validation.py:validate\_case}。
准入阈值统一为 $10^{-7}$；表中残差均已恢复为相应工程量，而非求解器内部缩放残差。

\subsection{实验 A：一时段解析最优解}

冻结输入为电负荷 $5$ MW、光伏可用 $2$ MW、购电价 $100$/MWh、光伏运维 $2$/MWh；其余载能、
设备、储能、碳与出口全部关闭。因光伏边际成本低于购电且容量受限，唯一经济调度为
\begin{equation}
 P^{pv}=2,qquad P^{imp}=5-2=3\ \text{MW},qquad
 C^*=2\times2+100\times3=304.
\end{equation}

\begin{center}\small
\begin{tabular}{@{}lrrr@{}}\toprule
量 & 解析值 & StrictHiGHS & 绝对误差\\\midrule
光伏 / MW & 2 & 2 & $0$\\
购电 / MW & 3 & 3 & $0$\\
目标/总成本 & 304 & 304 & $0$\\\bottomrule
\end{tabular}
\end{center}

该算例独立验证方向与目标系数，但不覆盖转换、库存或 MILP 分支。

\subsection{实验 B：两时段储能套利解析解}

两时段负荷均为 $1$ MW，购电价为 $(20,120)$/MWh；储能容量 $1$ MWh，初末 SOC 为零，
$\eta^{ch}=\eta^{dis}=0.9$、$\rho=1$，吞吐费与出口为零。低价时充满需要
$P_0^{ch}=1/0.9=1.111111111$ MW，高价时最多放出
$P_1^{dis}=0.9$ MW。故
\begin{equation}
 P^{imp}=(2.111111111,0.1),\quad E=(0,1,0),\quad C^*=54.2222222222.
\end{equation}

\begin{center}\small
\begin{tabular}{@{}lrrr@{}}\toprule
量 & 解析值 & StrictHiGHS & 最大绝对误差\\\midrule
充/放与中间 SOC & $(1.111111,0.9,1)$ & 同左 & $1.110\times10^{-16}$\\
目标/总成本 & $54.2222222222$ & $54.2222222222$ & $7.105\times10^{-15}$\\\bottomrule
\end{tabular}
\end{center}
无储能成本为 $140$，因此该最优解节省 $85.7777777778$。这项比较使用解析 LP，而非另一个调用同一 AML
模型的求解器包装。

\subsection{实验 C：24 时段全载能样例}

样例由 \srcpath{src/integrated\_energy/integrated\_energy\_optimizer.cpp:make\_sample\_campus\_ies\_data}
确定性生成，包含四载能、五类库存、交通、CCUS 与成本目标。StrictHiGHS 返回 optimal：总成本
$14039.18038252814$，购电 $52.6302726071$ MWh，可再生出力 $110.6860329018$ MWh，总排放
$68.7957633345$ tCO$_2$，残余碳 $12.7846107510$ tCO$_2$。

Python oracle 从输出向量独立重建每条作者方程，最大误差如下。
\begin{center}\small
\begin{tabular}{@{}lrlr@{}}\toprule
检查族 & 最大误差 & 检查族 & 最大误差\\\midrule
电平衡 / MW & $1.776\times10^{-15}$ & 热平衡 / MW & $4.441\times10^{-16}$\\
氢平衡 / MW & $4.441\times10^{-16}$ & 燃料平衡 / MW & $5.343\times10^{-16}$\\
五类库存 / MWh & $1.066\times10^{-14}$ & 转换等式 & $0$\\
可再生可用量 / MW & $9.797\times10^{-16}$ & 交通映射 & $0$\\
碳关系 / tCO$_2$ & $2.498\times10^{-15}$ & 总成本重算 & $1.819\times10^{-12}$\\\bottomrule
\end{tabular}
\end{center}

\subsection{证据强度与未覆盖项}

实验 A/B 是独立解析最优值比较；实验 C 是独立方程残差与目标重算，能发现符号、索引、单位和抽取错误，
但不能独立证明 24 h 解的全局最优性。三项都没有验证 AC 电压、无功或支路约束；也没有第二个外部 MILP
后端的全模型最优值对照。MIPSolvers 本地 HEAD 与仓库 pin 不一致，因此这些数字是当前工作树证据，
不是已固定依赖的发布基线。

\begin{gapnote}
要把 24 h 最优性升级为跨引擎证据，需要维护一个不复用 AML builder 的独立 LP/MILP 公式导出或外部
模型，并比较目标、离散模式和全部原始残差。当前证据只在上述三层范围内准入。
\end{gapnote}

\section{验证证据、独立复算与工程准入}

\subsection{注册测试覆盖}

当前唯一模块专用注册目标为 \srcpath{tests/test\_integrated\_energy\_campus.cpp}，由
\srcpath{tests/CMakeLists.txt} 注册为 \texttt{test\_integrated\_energy\_campus}。其六类用例覆盖：
\begin{enumerate}
  \item 24 步样例求解、四载能逐时平衡、状态长度和 Sankey carrier；
  \item Cost/Carbon/MinCurtailment/Weighted 目标切换及求解状态；
  \item 分层储氢开关、库存界限、转移变量和循环末态；
  \item 聚合 PCC、模型范围、有效性 flags 和限制字符串；
  \item 交通关闭时总里程为零，以及开启时 EV/HV/ICV 里程和为需求；
  \item 非法被动效率与非 unity 功率因数的拒绝路径。
\end{enumerate}
测试证明的是构造算例上的实现一致性，不是外部综合能源系统工具的交叉验证。

\subsection{结果独立复算}

任何纳入研究结果的原始解至少应执行以下后验残差：
\begin{align}
r_e&=\max_t|\text{electric supply}_t-\text{electric demand}_t|,\\
r_q&=\max_t|\text{heat supply}_t-\text{heat demand}_t|,\\
r_h&=\max_t|\text{hydrogen supply}_t-\text{hydrogen demand}_t|,\\
r_f&=\max_t|\text{fuel supply}_t-\text{fuel demand}_t|,\\
r_s&=\max_{s,t}|E_{s,t+1}-\widehat E_{s,t+1}|.
\end{align}
$\widehat E$ 必须使用第 3 章的源码保持率/效率，而不是调用者另一个储能约定。
同时复算可用新能源拆分、CHP 三个包络、CCUS 上限、碳预算、可再生比例和所有变量界限。

\subsection{准入矩阵}

\begin{center}\footnotesize
\begin{longtable}{@{}L{4.0cm}L{5.4cm}L{5.8cm}@{}}
\toprule
用途 & 最低证据 & 禁止的外推\\\midrule
\endfirsthead
\toprule 用途 & 最低证据 & 禁止的外推\\\midrule\endhead
教学/接口演示 & 注册测试、\fld{feasible}、平衡残差 & 不声称电网或设备安全\\
成本/碳情景比较 & 实际 solver/status/gap、同一输入边界、目标标度审计 & 不把非 optimal 解排序当全局最优排序\\
储能规划研究 & 末态口径、库存残差、步长/保持率解释、敏感性 & 不把短时域循环解外推季节调度\\
电网协同运行 & 另行 PF/OPF 回放 PCC 轨迹及网络限值 & 本模块单独结果不得准入\\
生产控制 & 独立状态估计、滚动更新、故障与不确定性闭环 & 当前离线模型不构成控制器\\
\bottomrule
\end{longtable}
\end{center}

\subsection{已知证据边界}

开发状态记录的最近聚焦基线为 6 个用例、106 个断言通过；该结果来自特定 Debug 配置且本地
MIPSolvers 状态与记录 pin 不一致。它不替代完整 Release CTest、sanitizer、另一 MILP 引擎逐变量交叉验证，
也不证明更长时域的性能。本文档修改后的实跑结果以
\srcpath{docs/overview/development\_status.md} 最新条目为准。

\subsection{后续闭环优先级}

首要缺口是独立的模型 residual evaluator 与至少两后端同一固定算例的目标/变量交叉验证；其次是把
数据清洗报告公开，避免非有限容量、COP 回退和比例重写不可见；电网耦合只能通过新增明确接口、
网络变量/约束和回放测试实现，不能通过更改 \fld{model\_scope} 文案宣称完成。


\section{跨模块工业级技术评价基线}

本章规定本手册所述模块进入工程算例、批量研究或生产服务前必须共同满足的评价基线。
具体模型公式、变量、约束和专用算法以正文相应章节为准；本章不扩大源码已经实现的范围。

\subsection{输入、输出、边界条件与数据结构审计}
输入审计应逐项检查有限性、单位、上下界、枚举值和引用完整性。系统对象必须先按所需层级完成
校验；AC 与 DC 母线采用域限定映射，组件稳定 \texttt{index}、向量位置和临时图索引不得混用。
输出审计必须记录单位、索引空间、模型范围、收敛或可行状态、后端、容差、迭代/节点计数和告警。
空系统、空时域、孤岛、重复 ID、零容量、退化约束与不可用可选后端均属于边界测试集。

\subsection{工程假设与数值稳定性分析}
模型使用的平衡、线性化、稳态、独立性、概率分布、时间离散或网络等值假设必须与应用场景匹配。
数值评价至少包含变量缩放、绝对/相对残差、可行性违反、条件数或枢轴质量、步长/阻尼、迭代停滞
和随机种子复现性。若采用正则化、平滑、回退、松弛或近似，应同时报告触发条件、参数值、对主指标
的敏感性以及是否改变模型口径；不得用“已返回结果”替代收敛或有效性证明。

\subsection{复杂度与资源分析}
令 $n_x$、$n_c$、$n_t$、$n_s$、$|V|$、$|E|$ 分别表示变量、约束、时段、场景、节点和边数。
报告应分别给出建模、求解、后处理的时间与内存。图遍历通常为 $O(|V|+|E|)$；逐时段/逐场景
流程至少线性增长为 $O(n_t n_s)$ 次子问题调用；稠密线性代数最坏可达 $O(n_x^3)$，稀疏方法则应
报告结构非零数、填充率和因子分解峰值内存；MILP/NLP 不给出虚假的多项式最坏界，而报告节点数、
间隙、时限和实测扩展曲线。复杂度结论必须基于固定硬件、固定后端和可复现算例族。

\subsection{异常工况处理}
非法输入应在进入求解器前返回结构化错误；不可行、无界、不收敛、时限、内存不足和后端缺失必须
相互区分。部分结果只有在对应 \fld{ValidityFlags}、\fld{model\_scope}、
\fld{model\_limitations} 或等价字段允许时才可使用。异常路径不得静默丢弃 AC/DC 资产、制造平衡
节点、把 NaN 改写为零或把近似解标为全局最优；系统原始对象应保持可恢复和可重试。

\subsection{与其他模块的接口关系}
接口链路遵循“工程语义模型 $\to$ 校验 $\to$ 投影/装配 $\to$ 求解或分析 $\to$ 结果回投影”。
跨模块传递时保留稳定 ID、单位、符号、时间基准和场景身份；消费潮流、OPF、动态或可靠性结果前，
先检查其收敛、覆盖范围和有效性标志。HTTP/JSON/Python 层只做语义保持适配，不重新解释求解结果。

\subsection{测试与验证建议}
最低验证矩阵包括：手算小例、标准算例、边界/异常输入、随机性质测试、算法或后端交叉校核、
序列化往返、稳定 ID 回投影、AC/DC 同号母线、重复运行确定性和规模扩展测试。数值算法还应加入
残差独立重算、雅可比有限差分或自动微分校核；近似模型应与高保真模型比较并给出误差分布，而非
只报告单个平均值。回归报告必须列明构建配置、算例版本、容差、通过/失败/跳过数及未验证范围。

\subsection{工业级评估指标}
\begin{itemize}
  \item \textbf{正确性}：守恒残差、约束最大违反、解析/基准误差、结果回投影一致性；
  \item \textbf{鲁棒性}：收敛率、不可行识别率、异常分类准确率、扰动敏感性与随机复现率；
  \item \textbf{性能}：端到端时延、吞吐量、峰值内存、并行效率与规模增长斜率；
  \item \textbf{可追溯性}：输入模式版本、稳定 ID 覆盖率、模型范围/限制字段完整率、告警可定位率；
  \item \textbf{工程有效性}：与标准、外部引擎或现场数据的偏差，以及误判/漏判对业务指标的影响。
\end{itemize}
指标应预先给出验收阈值和失败处置；未达到阈值时只能标记为实验性或受限适用。

\subsection{适用范围、局限性与后续改进}
本手册只评价当前源码覆盖的模型、后端和数据结构，不对未建模物理过程、未安装求解器或未执行的
外部验证作保证。后续改进应按“缺口 $\to$ 可量化指标 $\to$ 源码/测试变更 $\to$ 独立验证”闭环，
优先消除会造成错误结论的语义缺口，其次提升数值鲁棒性与性能，最后扩展模型范围。


\end{document}
